\section{Module Introduction}

Programming Paradigms has the module code COMP2221 and is led by \hyperref[modulelead-stuartjames]{Dr Stuart James}.

\subsection{Assessment}

This module is assessed through one piece of summative coursework and two final exams, these can be seen in \elemref{tab:Assessments for Programming Paradigms.}{Table}. \comment{Final dates to be confirmed.}

\inserttable{c|c|c|c|c|c}{Assessments for Programming Paradigms.}
{\textbf{Title} & \textbf{Opening} & \textbf{Due} & \textbf{Feedback} & \textbf{Professor} & \textbf{Weight}}
{
    Systems Programming & 31/10/2026 & 11/12/2026 & 08/01/2027 & \hyperref[modulelead-stuartjames]{Dr Stuart James} & $50\%$ \\
    Object Oriented $\&$ Functional & May Exam & n/a & n/a & \hyperref[lecturer-mladenivkovic]{Ivkovic} $\&$ \hyperref[lecturer-billymoses]{Moses} & $50\%$}

\subsection{Timetable}

This module is taught through two one-hour lectures and one two-hour computer class each week as seen below in \elemref{tab:Michaelmas Term Timetable for Programming Paradigms.}{Table}.

\inserttable{c|c|c|c|c}{Michaelmas Term Timetable for Programming Paradigms.}
{\textbf{Day} & \textbf{Time} & \textbf{Class Type} & \textbf{Location} & \textbf{Lecturer}}
{
    Monday & 11:00 - 12:00 & \colouredText{teal}{Lecture} & CLC202 & \hyperref[modulelead-stuartjames]{Dr Stuart James} \\
    Monday & 16:00 - 18:00 & \colouredText{blue}{Computer Class} & CB-LG001 & n/a \\
    Friday & 14:00 - 15:00 & \colouredText{teal}{Lecture} & CLC202 & \hyperref[modulelead-stuartjames]{Dr Stuart James}}

\section{Systems Programming}

This section is taught by \hyperref[modulelead-stuartjames]{Dr Stuart James}. There is no official reading list for this module.

\separate{5}\definition{Systems Programming}{\underline{Systems programming} involves the development of individual pieces of software that allow the entire system to function as a single unit.}

\subsection{Version Control and UNIX}

\subsection{The Basics of C}

\subsubsection{Basic File Setup}

\definition{pre-processor directives}{\underline{Pre-processor directives} are lines which start with a \# and are commands to the C pre-processor.
\begin{itemize}
    \item If the directive is \underline{\#include}, and \(<>\) are used, then \(/usr/include\) is prioritised in UNIX by convention, whereas if \say{} are used, then the current directory is prioritised.
    \item \definition{Definition Macro}{If the directive is \#define, then this is a \underline{definition macro}. These are used to provide definitions in code.}
    \item \definition{Conditionals}{You can also use directives for debugging without generating overheads, these are called \underline{conditionals}. They are as shown in \elemref{tab:Different types of conditionals.}{Table}.}
\end{itemize}}

\inserttable{c|c}{Different types of conditionals.}{\textbf{Conditional} & \textbf{Function}}
{\#ifdef VAR & Tests if VAR is \#defined. \\
 \#ifndef VAR & Tests if VAR is not \#defined. \\
 \#else & Runs if the original condition evaluates as false. \\
 \#endif & Ends the if testing condition.}

\noindent You can also put debugging code in the preamble, as seen in \prealgoref{alg:Debugging Code Example.}.

\noindent\hrulefill
\begin{minted}{c}
#define MY_DEBUG // Define an identifier.

#ifdef MY_DEBUG
    assert(i > 0);
    printf("i is %d \n", i);
#endif
\end{minted}
\hrule
\algorithmlabel{Debugging Code Example.}

\noindent \definition{Parameterised macro definitions}{You can also use \#define to define \underline{parameterised macros}. These act like functions, and an example can be seen in \prealgoref{alg:Parameterised Macro Example.}.}

\noindent\hrulefill
\begin{minted}{c}
#define ADD(a,b) a + b
\end{minted}
\hrule
\algorithmlabel{Parameterised Macro Example.}

\separate{7}\definition{the main function}{The \underline{main() function} is the entry function for C programs. You can only have one of these.}

\subsubsection{Printf Parameters}

\inserttable{c|c}{First parameter of printf.}{\textbf{Value} & \textbf{Resulting Format}}
{\(\%d\) & Signed decimal (int). \\
 \(\%u\) & Unsigned decimal. \\
 \(\%o,\%x\) & Octal and hexidecimal. \\
 \(\%l\) & Long. \\
 \(\%f\) & Floating point, so \(\%1.2f\) will give \(3.14\). \\
 \(\%e\) & Floating point (exponent form). \\
 \(\%c,\%s\) & Character and string.}

\subsection{Control Flow and Functions}

\subsection{Data Types and Structures}

\subsection{Advanced Bash Scripting}

\subsection{Pointers}

\subsection{Memory Management}

\subsection{Debugging}

\subsection{Scope and Recursion}

\subsection{Libraries}

\subsection{stdlib}

\subsection{C++}

\section{Functional Programming}

This section is taught by \hyperref[lecturer-billymoses]{Dr Billy Moses Jr.}. \comment{Check for recommended reading.}

\section{Object Oriented Programming}

This section is taught by \hyperref[lecturer-mladenivkovic]{Dr Mladen Ivkovic}. \comment{Check for recommended reading.}

\subsection{Java Fundamentals}

\subsection{OOP with Classes and Objects}

\subsection{Encapsulation}

\subsection{Inheritance}

\subsection{Abstraction with Abstract Classes}

\subsection{Abstraction with Interfaces}

\subsection{Polymorphism}

\subsection{Java Collections}

\subsection{Best Practises}